% !TEX root = userguide.tex

\def\headdate{4th April 2017}

\section{Surface convection diffusion}
\label{sec:surface_convection_diffusion}

(These examples have been based on the codes developed by Christos Mitrias of the TU/e)

\subsection{Surfactant concentration on the interface of a drop subject to simple shear}

Example \texttt{drop12.f90} involves two Newtonian fluids, separated by an interface with 
surface tension which depends linearly on the surfactant concentration. 
The concentration of the surfactants on the interface is calculated by solving the following  
convection-diffusion equation on the interface:
\begin{equation}\label{eq:surfactant_concentration}
	\frac{\partial c}{\partial t} +  
	(\vec u-\vec u_\text{m})\cdot \nabla_\text{s} c + (\nabla_\text{s} \cdot \vec u )c = 
	\nabla_\text{s} \cdot ( D \nabla_\text{s} c ),
\end{equation}
where $D$ is the diffusion coefficient, $c$ is the surfactant concentration and $\vec u_\text{m}$ is the grid velocity on the interface.

The surfactant concentration equation is coupled with the flow via the velocity 
on the interface and the surface tension which is given by
\begin{equation}
	\Gamma(c) = \frac{\Gamma_0}{1-\beta} (1-\beta \frac{c}{c_0}),
\end{equation}
where $\Gamma_0$ is the surface tension at the reference concentration $c_0$ and $\beta$ is a constant.

In Figure~\ref{fig:surfactant_concentration} the result of the surfactant concentration 
on a drop interface that has reached steady state is shown.
\begin{figure}[htbp!]
   \begin{center}
   \includegraphics[scale=0.4]{surfactant_concentration}
   \end{center}
   \caption{Steady state of a drop in simple shear containg insoluble surfactants on 
	   the interface. Results obtained using example
   \texttt{drop12.f90}}
    \label{fig:surfactant_concentration}
\end{figure}

\subsection{Surfactant concentration on the interface of a drop subject to uniaxial extensional flow (axisymmetric)}

Example \texttt{drop14.f90} involves two Newtonian fluids, separated by an interface with 
surface tension which depends linearly on the surfactant concentration. Axisymmetry is assumed. 
The example is based on the the example \texttt{drop13.f90} (see Sec.~\ref{sec:drop_remesh_axisymmetric}).
